% Options for packages loaded elsewhere
\PassOptionsToPackage{unicode}{hyperref}
\PassOptionsToPackage{hyphens}{url}
\PassOptionsToPackage{dvipsnames,svgnames,x11names}{xcolor}
%
\documentclass[
  11pt,
]{article}
\usepackage{amsmath,amssymb}
\usepackage{iftex}
\ifPDFTeX
  \usepackage[T1]{fontenc}
  \usepackage[utf8]{inputenc}
  \usepackage{textcomp} % provide euro and other symbols
\else % if luatex or xetex
  \usepackage{unicode-math} % this also loads fontspec
  \defaultfontfeatures{Scale=MatchLowercase}
  \defaultfontfeatures[\rmfamily]{Ligatures=TeX,Scale=1}
\fi
\usepackage{lmodern}
\ifPDFTeX\else
  % xetex/luatex font selection
  \setmainfont[]{DejaVu Serif}
  \setmonofont[Scale=0.82]{DejaVu Sans Mono}
\fi
% Use upquote if available, for straight quotes in verbatim environments
\IfFileExists{upquote.sty}{\usepackage{upquote}}{}
\IfFileExists{microtype.sty}{% use microtype if available
  \usepackage[]{microtype}
  \UseMicrotypeSet[protrusion]{basicmath} % disable protrusion for tt fonts
}{}
\makeatletter
\@ifundefined{KOMAClassName}{% if non-KOMA class
  \IfFileExists{parskip.sty}{%
    \usepackage{parskip}
  }{% else
    \setlength{\parindent}{0pt}
    \setlength{\parskip}{6pt plus 2pt minus 1pt}}
}{% if KOMA class
  \KOMAoptions{parskip=half}}
\makeatother
\usepackage{xcolor}
\usepackage[top=2.5cm,bottom=2.5cm,left=2.8cm,right=2.8cm]{geometry}
\usepackage{listings}
\newcommand{\passthrough}[1]{#1}
\lstset{defaultdialect=[5.3]Lua}
\lstset{defaultdialect=[x86masm]Assembler}
\setlength{\emergencystretch}{3em} % prevent overfull lines
\providecommand{\tightlist}{%
  \setlength{\itemsep}{0pt}\setlength{\parskip}{0pt}}
\setcounter{secnumdepth}{-\maxdimen} % remove section numbering
\ifLuaTeX
  \usepackage{selnolig}  % disable illegal ligatures
\fi
\IfFileExists{bookmark.sty}{\usepackage{bookmark}}{\usepackage{hyperref}}
\IfFileExists{xurl.sty}{\usepackage{xurl}}{} % add URL line breaks if available
\urlstyle{same}
\hypersetup{
  colorlinks=true,
  linkcolor={Maroon},
  filecolor={Maroon},
  citecolor={Blue},
  urlcolor={Blue},
  pdfcreator={LaTeX via pandoc}}

\title{SpeakPrep ��� Document 6: Agent Documentation Prompts}
\author{Abhiyan Sainju}
\date{April 2026}

\begin{document}
\maketitle

\hypertarget{speakprep-document-6-agent-documentation-prompts}{%
\section{SpeakPrep --- Document 6: Agent Documentation
Prompts}\label{speakprep-document-6-agent-documentation-prompts}}

\hypertarget{auto-generate-your-build-log-learnings-decisions-blog-posts-linkedin-and-more}{%
\subsubsection{Auto-generate your build log, learnings, decisions, blog
posts, LinkedIn, and
more}\label{auto-generate-your-build-log-learnings-decisions-blog-posts-linkedin-and-more}}

\hypertarget{version-1.0-abhiyan-sainju-april-2026}{%
\subsubsection{Version 1.0 \textbar{} Abhiyan Sainju \textbar{} April
2026}\label{version-1.0-abhiyan-sainju-april-2026}}

\begin{center}\rule{0.5\linewidth}{0.5pt}\end{center}

\begin{quote}
These are prompts you give to an agent (Claude, Cursor, Claude Code) to
write your documentation FOR you. You provide the raw input (what you
did, what confused you, what you decided). The agent writes it up
properly. Run each prompt at the right time. The agent is your
ghostwriter --- you supply the experience, it supplies the words.
\end{quote}

\begin{center}\rule{0.5\linewidth}{0.5pt}\end{center}

\hypertarget{how-this-works}{%
\subsection{HOW THIS WORKS}\label{how-this-works}}

You never write documentation from scratch. Instead:

\begin{enumerate}
\def\labelenumi{\arabic{enumi}.}
\tightlist
\item
  At the end of every dev session: run the \textbf{Daily Build Log
  Prompt}
\item
  When something clicks or confuses you: run the \textbf{Learnings
  Prompt}
\item
  When you make a technical choice: run the \textbf{Decision Prompt}
\item
  At the end of each phase: run the \textbf{Phase Summary Prompt}
\item
  When you want to post on LinkedIn: run the \textbf{LinkedIn Prompt}
\item
  When you want to write a blog post: run the \textbf{Blog Post Prompt}
\item
  When you want a deep technical write-up: run the \textbf{Technical
  Deep Dive Prompt}
\item
  At the end of the whole project: run the \textbf{Full Journey Prompt}
\end{enumerate}

\textbf{Which agent to use for all of these:} Claude (this chat) or
Claude Code. Not Cursor --- Cursor is for code editing. These are
writing tasks.

\textbf{The input you give the agent:} Just talk naturally. Voice notes
transcribed, bullet points, messy notes --- anything. The agent's job is
to turn your raw thoughts into clean documentation. You don't need to
write well. You just need to remember what happened.

\begin{center}\rule{0.5\linewidth}{0.5pt}\end{center}

\hypertarget{file-structure}{%
\subsection{FILE STRUCTURE}\label{file-structure}}

\begin{lstlisting}
speakprep/
├── docs/
│   ├── journey/
│   │   ├── BUILD_LOG.md          ← daily raw entries, one per session
│   │   ├── DECISIONS.md          ← every technical/product decision
│   │   ├── LEARNINGS.md          ← concepts explained in your own words
│   │   ├── BUGS.md               ← bugs you hit and how you fixed them
│   │   └── phase-summaries/
│   │       ├── phase0-summary.md
│   │       ├── phase1-summary.md
│   │       └── ...
│   ├── content/
│   │   ├── linkedin-posts/       ← ready-to-post LinkedIn content
│   │   ├── blog-posts/           ← full blog articles
│   │   ├── technical-deepdives/  ← detailed technical write-ups
│   │   └── tweet-threads/        ← Twitter/X thread drafts
│   └── final/
│       ├── full-journey.md       ← complete story from start to finish
│       └── research-notes.md     ← for a potential paper or talk
\end{lstlisting}

Create this structure now:

\begin{lstlisting}[language=bash]
mkdir -p docs/journey/phase-summaries
mkdir -p docs/content/linkedin-posts
mkdir -p docs/content/blog-posts
mkdir -p docs/content/technical-deepdives
mkdir -p docs/content/tweet-threads
mkdir -p docs/final
touch docs/journey/BUILD_LOG.md
touch docs/journey/DECISIONS.md
touch docs/journey/LEARNINGS.md
touch docs/journey/BUGS.md
\end{lstlisting}

\begin{center}\rule{0.5\linewidth}{0.5pt}\end{center}

\hypertarget{prompt-1-daily-build-log}{%
\subsection{PROMPT 1: DAILY BUILD LOG}\label{prompt-1-daily-build-log}}

\textbf{When to run:} At the end of every dev session. Even short ones.
\textbf{Time it takes you:} 5 minutes of talking/typing your raw notes.
\textbf{What the agent produces:} A clean, dated BUILD\_LOG.md entry.

\textbf{How to give input:} Just dump everything into the prompt below.
Don't organize it. Write like you're texting a friend about what you
did. Messy is fine.

\begin{center}\rule{0.5\linewidth}{0.5pt}\end{center}

\hypertarget{the-prompt}{%
\subsubsection{The Prompt:}\label{the-prompt}}

\begin{lstlisting}
You are helping me document my development journey building SpeakPrep,
a real-time voice AI mock interview coach. I am learning while building —
this documentation will later become blog posts, LinkedIn content, and
a full journey write-up.

Today's raw notes (my stream of consciousness — do NOT clean up my meaning,
just clean up the writing):
---
[PASTE YOUR RAW NOTES HERE — everything you did, tried, broke, fixed,
learned, felt confused about, felt proud of. Voice note transcript is fine.]
---

Write a BUILD_LOG.md entry for today with this structure:

## [Date] — Day [N] — Phase [X]: [one-line summary of the session]

### What I Built Today
[2-4 sentences: what actually got done. Concrete. Specific file names, function names.]

### What I Learned
[Bullet list: each bullet is one concept or insight. Write it like explaining
to a smart friend who isn't a developer. Simple language, real understanding.]

### What Confused Me (and how I resolved it, or if still open)
[Honest account of what was hard. Don't skip this — it's the most valuable
part for future readers who will feel the same confusion.]

### Decisions Made
[Any choices between approaches. Why you picked one over another.
Even small ones: "used X instead of Y because..."]

### Bugs / Errors Hit
[Any error messages, unexpected behavior, things that didn't work at first.
What you tried. What finally fixed it.]

### How It Feels
[One honest sentence about energy, momentum, confidence. This is for the
human story in the blog/LinkedIn content later.]

### Tomorrow's Plan
[1-3 specific tasks, not vague goals]

After writing the entry, also:
1. Extract any new DECISIONS and format them for DECISIONS.md
   (format: ### D-[N] | [date] | [title] / Why / Tradeoff / What I'd do differently)
2. Extract any new LEARNINGS and format them for LEARNINGS.md
   (format: ### [concept] — [date] / [explanation in your own words, 3-8 sentences])
3. Extract any BUGS worth documenting for BUGS.md
   (format: ### BUG-[N] | [date] | [title] / Symptom / Root cause / Fix / Lesson)

Output each section clearly labeled so I can copy into the right file.
\end{lstlisting}

\begin{center}\rule{0.5\linewidth}{0.5pt}\end{center}

\hypertarget{prompt-2-learnings-entry-on-demand}{%
\subsection{PROMPT 2: LEARNINGS ENTRY
(on-demand)}\label{prompt-2-learnings-entry-on-demand}}

\textbf{When to run:} The moment something clicks. When you finally
understand why asyncio works the way it does, why 400ms is the right
silence threshold, why JWT tokens expire in 15 minutes --- whatever it
is. Run this immediately while it's fresh. Don't wait until end of day.

\textbf{How to give input:} Describe the concept and what made it click.
Can be one sentence: ``I just understood why you can't use time.sleep()
in async code.'' That's enough.

\begin{center}\rule{0.5\linewidth}{0.5pt}\end{center}

\hypertarget{the-prompt-1}{%
\subsubsection{The Prompt:}\label{the-prompt-1}}

\begin{lstlisting}
I just had a learning moment while building SpeakPrep. Help me document it
properly in LEARNINGS.md so I can reference it later and use it in blog posts.

What I just understood:
---
[DESCRIBE WHAT CLICKED — can be messy, one sentence, or a paragraph.
What was confusing before? What's clear now? What's the mental model?
What would you tell past-you to understand it faster?]
---

Context (optional — fill in what's relevant):
- Phase I'm in: [Phase X]
- Task I was doing: [task name]
- The specific thing that made it click: [code, explanation, experiment]

Write a LEARNINGS.md entry in this format:

### [Concept Name] — [Date]

**The confusion before:**
[What didn't make sense. What I thought it was. Why that was wrong.]

**The mental model that works:**
[The explanation in simple terms. Use an analogy if there is a good one.
Write it like explaining to a smart non-developer friend.]

**Why it matters for SpeakPrep specifically:**
[How does this concept actually show up in what I'm building?
What would break if I misunderstood it?]

**The code that made it real:**
[2-5 lines of the specific code that illustrated the concept.
Not the whole function — just the key line(s).]

**What to read if you want to go deeper:**
[One specific resource — article title + URL, or book + chapter.]

**Interview answer version:**
[How would I explain this concept in a technical interview in 3 sentences?
This is for when I'm asked "explain asyncio" or "how does VAD work."]
\end{lstlisting}

\begin{center}\rule{0.5\linewidth}{0.5pt}\end{center}

\hypertarget{prompt-3-decision-entry-on-demand}{%
\subsection{PROMPT 3: DECISION ENTRY
(on-demand)}\label{prompt-3-decision-entry-on-demand}}

\textbf{When to run:} Every time you choose between two technical
approaches. Even small ones. ``Should I use X or Y?'' --- log it when
you decide.

\begin{center}\rule{0.5\linewidth}{0.5pt}\end{center}

\hypertarget{the-prompt-2}{%
\subsubsection{The Prompt:}\label{the-prompt-2}}

\begin{lstlisting}
I just made a technical decision while building SpeakPrep. Help me document
it properly in DECISIONS.md for future reference and blog content.

The decision:
---
[DESCRIBE THE CHOICE YOU MADE AND WHY — can be a few sentences.
What were the options? What did you pick? What made you pick it?]
---

Write a DECISIONS.md entry in this format:

### D-[N] | [Date] | [Short title of the decision]

**The question I was facing:**
[What problem needed a decision? What were the stakes?]

**Options I considered:**
Option A — [name]: [what it is, pros, cons]
Option B — [name]: [what it is, pros, cons]
(Option C if applicable)

**What I chose:** Option [X]

**Why:**
[Real reasoning. Not just "it's better" — what specifically about my
situation made this the right call? What would have to be different
for me to pick the other option?]

**What I'm giving up:**
[Honest tradeoff. Every decision has a cost. What's mine?]

**How I'll know if I was wrong:**
[What signal would tell me this was the wrong choice?
What would make me switch?]

**Interview answer version:**
[How would I explain this decision in a system design interview
or "tell me about a technical decision you made" question? 2-3 sentences.]
\end{lstlisting}

\begin{center}\rule{0.5\linewidth}{0.5pt}\end{center}

\hypertarget{prompt-4-bug-story-on-demand}{%
\subsection{PROMPT 4: BUG STORY
(on-demand)}\label{prompt-4-bug-story-on-demand}}

\textbf{When to run:} After fixing a bug that took more than 15 minutes
to solve, or any bug that taught you something important.

\begin{center}\rule{0.5\linewidth}{0.5pt}\end{center}

\hypertarget{the-prompt-3}{%
\subsubsection{The Prompt:}\label{the-prompt-3}}

\begin{lstlisting}
I just fixed a bug while building SpeakPrep. Help me document it properly
in BUGS.md. This will become one of my "bug story" interview answers later.

The bug:
---
[DESCRIBE EVERYTHING — what the symptom was, what you tried, what you
thought it was at first, what it actually was, how you fixed it.
Include any error messages you saw.]
---

Write a BUGS.md entry AND an interview-ready bug story:

### BUG-[N] | [Date] | [Short title]

**Phase:** [Phase X — Task Y]
**Time to fix:** [rough estimate]

**The symptom:**
[What I observed. What the error message said. What behavior I expected
vs what actually happened.]

**My debugging process:**
[Step by step what I tried. What led me in the wrong direction.
What finally pointed me to the real cause. Be specific — this is the
educational part.]

**Root cause:**
[The actual reason it was happening. One clear sentence.]

**The fix:**
[What I changed. Key code lines if relevant.]

**Why it happened:**
[The underlying concept I misunderstood or didn't know.
What would have prevented this bug if I'd known it earlier?]

**The lesson:**
[One sentence: "Now I always [do X]" or "Never [do Y] because [Z]."]

---
Then write this as an interview STAR story:

**Interview version (STAR format):**
Situation: [context — what I was building, what phase]
Task: [what I needed to make work]
Action: [how I debugged and fixed it — specific steps]
Result: [what got fixed, what I learned, what I'd tell a junior dev]
\end{lstlisting}

\begin{center}\rule{0.5\linewidth}{0.5pt}\end{center}

\hypertarget{prompt-5-phase-summary}{%
\subsection{PROMPT 5: PHASE SUMMARY}\label{prompt-5-phase-summary}}

\textbf{When to run:} When you complete a phase (Phase 0, 1, 2, etc.).
\textbf{Input:} Your BUILD\_LOG.md entries from that phase + your own
reflections.

\begin{center}\rule{0.5\linewidth}{0.5pt}\end{center}

\hypertarget{the-prompt-4}{%
\subsubsection{The Prompt:}\label{the-prompt-4}}

\begin{lstlisting}
I just completed [Phase X — Phase Name] of building SpeakPrep.
Help me write a phase summary document.

My BUILD_LOG entries from this phase:
---
[PASTE ALL BUILD_LOG.md ENTRIES FROM THIS PHASE]
---

My own reflections (messy notes):
---
[ANYTHING YOU WANT TO ADD — how it felt, what surprised you,
what took longer than expected, what you're proud of]
---

Write a phase summary at docs/journey/phase-summaries/phase[X]-summary.md:

# Phase [X] — [Phase Name]: Complete
## [Start Date] → [End Date] | [Total days]

### What This Phase Built
[2-3 sentences: the concrete output. What exists now that didn't before?]

### The Most Important Thing I Learned
[One concept from this phase that changed how I think about building software.
Not a list — pick the one that mattered most and explain it well.]

### The Hardest Moment
[The specific point where I was most stuck or confused.
What it felt like. How I got through it.]

### The Proudest Moment
[The specific moment something worked that hadn't before.
What it felt like when it clicked.]

### By The Numbers
- Days taken: [N]
- Commits made: [run: git log --oneline | wc -l for this branch]
- Files created: [N]
- Lines of code written: [rough estimate]
- Bugs hit and fixed: [N]
- Times I wanted to quit: [honest number]

### What I'd Do Differently
[If I started this phase again tomorrow, knowing what I know now,
what would I do in the first hour that I didn't do?]

### What's Next
[Phase X+1 preview: what I'm about to build, what I'm nervous about]

---
Then write a SHORT version (3 paragraphs) suitable for a LinkedIn update post:
Paragraph 1: What I built this phase (concrete, specific)
Paragraph 2: The one thing that surprised me
Paragraph 3: What's next
End with 3-5 relevant hashtags.
\end{lstlisting}

\begin{center}\rule{0.5\linewidth}{0.5pt}\end{center}

\hypertarget{prompt-6-linkedin-post-on-demand}{%
\subsection{PROMPT 6: LINKEDIN POST
(on-demand)}\label{prompt-6-linkedin-post-on-demand}}

\textbf{When to run:} Any time you want to post. Can be mid-phase.
\textbf{Input:} A specific thing you built, learned, or realized. One
thing per post.

\begin{center}\rule{0.5\linewidth}{0.5pt}\end{center}

\hypertarget{the-prompt-5}{%
\subsubsection{The Prompt:}\label{the-prompt-5}}

\begin{lstlisting}
Help me write a LinkedIn post about something I did while building SpeakPrep.

What I want to post about:
---
[DESCRIBE THE SPECIFIC THING — could be a technical insight, a bug you fixed,
a milestone you hit, something that surprised you, a decision you made.
One thing. Be specific. Don't try to summarize everything.]
---

My tone and audience: I'm a recent MS CS grad from GW actively job hunting.
My audience is software engineers, hiring managers, and recruiters.
I want to come across as: genuinely learning, technically competent, honest
about challenges, not performatively humble or braggy.

Write 3 versions of the LinkedIn post:

VERSION A — Technical insight (for engineers):
Lead with the technical thing you learned. Explain it simply. Show the code
or the concept. End with why it matters. 150-200 words. No fluff.

VERSION B — Journey/story (for hiring managers):
Lead with the human moment — the confusion, the breakthrough, the decision.
Connect it to building something real. End with what you'd do differently.
150-200 words. Personal but professional.

VERSION C — Short punchy (for everyone):
3-5 lines max. Hook → specific insight → implication. No bullet points.
The kind of post that gets read in 10 seconds.

Rules for all versions:
- No "excited to share" or "humbled to announce"
- No excessive exclamation marks
- First sentence must make someone stop scrolling
- Include the project name (SpeakPrep) naturally
- End with a question or observation that invites comments
- 3-5 hashtags at the end (not inline)
\end{lstlisting}

\begin{center}\rule{0.5\linewidth}{0.5pt}\end{center}

\hypertarget{prompt-7-blog-post}{%
\subsection{PROMPT 7: BLOG POST}\label{prompt-7-blog-post}}

\textbf{When to run:} After completing a phase, or when you have a
specific technical topic you want to write about deeply. \textbf{Best
topics:} One concept per post. ``How I built X'' or ``Why I chose X over
Y'' or ``What I learned about X building a voice AI.''

\begin{center}\rule{0.5\linewidth}{0.5pt}\end{center}

\hypertarget{the-prompt-6}{%
\subsubsection{The Prompt:}\label{the-prompt-6}}

\begin{lstlisting}
Help me write a technical blog post for my developer blog / portfolio.

Topic: [ONE SPECIFIC TOPIC — e.g., "How streaming TTS works in a voice AI pipeline"
or "Why I chose Groq over OpenAI for real-time voice" or "What I learned
about Python asyncio building a WebSocket server"]

My raw notes and experience on this topic:
---
[PASTE YOUR LEARNINGS.md ENTRY + BUILD_LOG entries + any notes on this topic]
---

Audience: Software engineers with 1-5 years experience. They know Python
and web development but may not know voice AI, asyncio deeply, or the
specific tools I used.

Write a complete blog post with:

# [Title — specific, searchable, not clickbait]

## Introduction (150 words)
Hook: the problem or question. Why does this matter?
What will the reader know by the end?

## [Section 1] — Background / The Problem (200-300 words)
What was I trying to do? What didn't work at first?
This is where the reader gets context.

## [Section 2] — The Concept Explained (300-400 words)
The core technical concept, explained simply.
Include a diagram (described in text — I'll draw it) if it helps.
Real code snippets from my actual project (short, annotated).

## [Section 3] — How I Implemented It (300-400 words)
The specific approach I took. Key code. What I tried first that didn't work.

## [Section 4] — What Surprised Me / What I'd Do Differently (150-200 words)
The honest take. What was harder than expected? What shortcut did I miss?

## Conclusion (100 words)
What to take away. One clear sentence that summarizes the lesson.
Link to the repo.

---
Rules:
- Code snippets must be real (from my actual project files)
- Every technical claim must be explained, not assumed
- Conversational but technically rigorous
- No padding — if a sentence doesn't add value, cut it
- Approximate total length: 1,200-1,800 words
- Suggest 3 SEO-friendly titles and I'll pick one
\end{lstlisting}

\begin{center}\rule{0.5\linewidth}{0.5pt}\end{center}

\hypertarget{prompt-8-technical-deep-dive}{%
\subsection{PROMPT 8: TECHNICAL DEEP
DIVE}\label{prompt-8-technical-deep-dive}}

\textbf{When to run:} When you want a thorough technical document on one
concept. Good for portfolio, for reference, or as the foundation of a
talk or paper.

\begin{center}\rule{0.5\linewidth}{0.5pt}\end{center}

\hypertarget{the-prompt-7}{%
\subsubsection{The Prompt:}\label{the-prompt-7}}

\begin{lstlisting}
Help me write a technical deep dive document on a concept I learned
while building SpeakPrep.

Topic: [SPECIFIC CONCEPT — e.g., "Real-time audio streaming pipeline architecture"
or "ELO-based adaptive difficulty in interview coaching systems"
or "Multi-provider LLM fallback with circuit breakers"]

My experience and notes:
---
[PASTE EVERYTHING RELEVANT — LEARNINGS.md entries, BUILD_LOG notes,
the actual code you wrote, decisions you made, bugs you hit]
---

Write a technical deep dive at docs/content/technical-deepdives/[topic].md

Structure:

# [Topic]: A Deep Dive from Building SpeakPrep

## Abstract (100 words)
What this covers. Who should read it. What they'll be able to do after.

## Problem Statement
What problem does this concept solve?
Why is the naive approach insufficient?

## Core Concepts
[Build up the theory from first principles.
Define every term before using it.
Use diagrams described in ASCII or text form.
Link to academic papers or official docs where relevant.]

## Implementation
[The actual code I wrote. Annotated heavily.
Show the evolution: what I tried first, why it didn't work,
what the final version looks like and why.]

## Performance Characteristics
[Numbers where I have them. Latency, throughput, memory.
What scales, what doesn't.]

## Tradeoffs and Alternatives
[What else I considered. Why I didn't use it.
When would you use the alternative instead?]

## Lessons Learned
[What I'd tell someone starting this from scratch.
What the docs don't tell you.
What only experience teaches.]

## References
[Papers, blog posts, documentation that helped me understand this.
With brief notes on what each one contributed.]

Target length: 2,000-4,000 words depending on complexity.
Write for a senior engineer audience — assume strong fundamentals.
\end{lstlisting}

\begin{center}\rule{0.5\linewidth}{0.5pt}\end{center}

\hypertarget{prompt-9-tweet-x-thread}{%
\subsection{PROMPT 9: TWEET / X THREAD}\label{prompt-9-tweet-x-thread}}

\textbf{When to run:} When you have a specific technical insight that
can be explained in a short thread. Good for building audience.

\begin{center}\rule{0.5\linewidth}{0.5pt}\end{center}

\hypertarget{the-prompt-8}{%
\subsubsection{The Prompt:}\label{the-prompt-8}}

\begin{lstlisting}
Help me write a Twitter/X thread about something I built or learned
while building SpeakPrep.

The insight:
---
[ONE SPECIFIC THING — a concept, a decision, a bug story, a surprising number.
Something concrete and specific. Not "here's everything I learned."]
---

Write a Twitter thread:

Tweet 1 (hook — max 280 chars):
Stop people scrolling. State the most surprising or counterintuitive thing.
Make them want to read the rest. Can be a question or a bold statement.

Tweet 2-3 (setup — max 280 chars each):
Build the context. What was I trying to do? What's the problem?

Tweet 4-6 (the insight — max 280 chars each):
The actual technical content. Be specific. Include code snippets where
they fit (use code blocks). Numbers > vague claims.

Tweet 7 (the lesson — max 280 chars):
One sentence takeaway. What should the reader remember?

Tweet 8 (CTA — max 280 chars):
Point to the repo, ask a question, or invite engineers who've faced
the same thing to share their approach.

Rules:
- Each tweet must stand alone (someone might screenshot one)
- No "🧵 Thread:" — just start with the hook
- Max 2 emojis per tweet, used sparingly
- No "in this thread I will..." — show don't announce
- Punchy. Short sentences. Real numbers.
\end{lstlisting}

\begin{center}\rule{0.5\linewidth}{0.5pt}\end{center}

\hypertarget{prompt-10-full-journey-document}{%
\subsection{PROMPT 10: FULL JOURNEY
DOCUMENT}\label{prompt-10-full-journey-document}}

\textbf{When to run:} Once. At the end of the project when it's
launched. \textbf{Input:} ALL your BUILD\_LOG.md entries + phase
summaries + your reflections. \textbf{Output:} The complete story of
building SpeakPrep, start to finish. This is the foundation for a
long-form blog post, a talk, or a research paper.

\begin{center}\rule{0.5\linewidth}{0.5pt}\end{center}

\hypertarget{the-prompt-9}{%
\subsubsection{The Prompt:}\label{the-prompt-9}}

\begin{lstlisting}
I have finished building SpeakPrep — a real-time voice AI mock interview coach
I built from scratch over [N months] while actively job hunting after graduating
with my MS in CS from George Washington University.

Here is all my documentation:

BUILD_LOG.md (all entries):
---
[PASTE FULL BUILD_LOG.MD]
---

Phase summaries:
---
[PASTE ALL PHASE SUMMARY FILES]
---

DECISIONS.md:
---
[PASTE DECISIONS.MD]
---

LEARNINGS.md:
---
[PASTE LEARNINGS.MD]
---

My personal context:
- I am Abhiyan Sainju, originally from Kathmandu, Nepal
- MS Computer Science, GW, December 2025, GPA 3.97
- Was actively job hunting while building this
- This project was motivated by my own frustration with interview prep tools
- I learned most of these technologies (asyncio, WebSockets, voice AI) while building

Write docs/final/full-journey.md — the complete story:

# Building SpeakPrep: A Developer's Journey
## From frustration to a real-time voice AI in [N months]

### Prologue: Why I Built This
[The personal motivation. The job search frustration. The gap I saw.
Write this as a story, not a product pitch. 300 words.]

### The Decision to Build
[What I decided to build and why. The key early decisions.
What I almost built instead. 200 words.]

### Phase by Phase: What Actually Happened
[For each phase, 300-400 words:
- What I planned to build
- What actually happened (different from the plan?)
- The hardest moment
- The breakthrough moment
- What I learned that I didn't expect to learn]

### The Technical Decisions That Mattered Most
[The 5 most important decisions from DECISIONS.md.
For each: what the choice was, why it mattered, would I make the same choice again?]

### What I Got Wrong
[Honest retrospective. 3-5 things that took longer or went differently
than expected. What I'd do differently. 300 words.]

### What Surprised Me Most
[The 3 most surprising things — technical, personal, or product. 200 words.]

### What I Know Now That I Didn't Know Then
[The most important things I learned — from LEARNINGS.md.
Not a list of technologies. The actual mental models and insights. 400 words.]

### The Numbers
[Concrete stats: days building, commits, bugs fixed, latency achieved,
users in beta, etc. Whatever you have.]

### What's Next
[For SpeakPrep, for my career, for what I want to build next. 200 words.]

### Epilogue
[One paragraph. What this project means to you.
What you'd tell someone starting this today. Personal, honest, short.]

---
Target length: 4,000-6,000 words.
Tone: Personal and honest. Not a press release. Not a tutorial.
The reader should feel like they were there watching it happen.
Write it in first person throughout.
After the full document, also produce:
1. A 500-word blog post version (the highlights only)
2. A 200-word LinkedIn post version
3. A 100-word bio blurb for talks/conferences: "Abhiyan built X while..."
\end{lstlisting}

\begin{center}\rule{0.5\linewidth}{0.5pt}\end{center}

\hypertarget{daily-habit-which-prompt-when}{%
\subsection{DAILY HABIT: WHICH PROMPT,
WHEN}\label{daily-habit-which-prompt-when}}

\begin{lstlisting}
Every dev session end (5 min):     → Prompt 1: Daily Build Log
Concept clicks mid-session (3 min): → Prompt 2: Learnings Entry
Make a technical choice (2 min):    → Prompt 3: Decision Entry
Fix a hard bug (5 min):            → Prompt 4: Bug Story
Finish a phase (20 min):           → Prompt 5: Phase Summary
Want to post on LinkedIn (10 min): → Prompt 6: LinkedIn Post
Want to write an article (30 min): → Prompt 7: Blog Post
Want deep technical doc (45 min):  → Prompt 8: Technical Deep Dive
Want a Twitter thread (10 min):    → Prompt 9: Tweet Thread
Project finished (2-3 hours):      → Prompt 10: Full Journey
\end{lstlisting}

\hypertarget{one-rule}{%
\subsection{ONE RULE}\label{one-rule}}

\textbf{Never give the agent a vague input.} ``I learned some stuff
about async today'' produces bad output. ``I spent 2 hours confused
about why time.sleep() inside an async function froze my entire server,
then realized it blocks the single event loop thread --- I tested it by
adding a print statement before and after and confirming nothing else
ran during the sleep'' produces great output.

Specific input → specific output. The agent can only work with what you
give it.

\begin{center}\rule{0.5\linewidth}{0.5pt}\end{center}

\emph{End of Document 6 --- Agent Documentation Prompts}

\end{document}
